\section{Methodology}
\label{sec:methodology}

\subsection{Research Questions and Study Scope}
\label{sec:research-questions}

We organize the study around three research questions
(Table~\ref{tab:research-questions}). The computational protocol treats a user
story and its acceptance criteria as a requirement case, with the system purpose
providing context. It assesses each catalog pattern against each case
independently. The assessment distinguishes applicability, non-applicability,
and insufficient information; a positive classification requires a supporting
user-story or acceptance-criterion reference. These classifications are model
outputs, not independently established ground truth.

\begin{table*}[t]
\caption{Research questions, units of analysis, and evidence mapping.}
\label{tab:research-questions}
\centering
\small
\begin{tabular}{p{0.05\textwidth} p{0.36\textwidth} p{0.49\textwidth}}
\hline
\textbf{ID} & \textbf{Research question} & \textbf{Unit, evidence, and intended analysis} \\
\hline
RQ1 &
How does the LLM classify architectural-pattern applicability across the requirement cases, and which patterns does it identify as applicable? &
Unit: one user-story--pattern pair.
Evidence: the official exhaustive assessment log, original requirements, and pattern catalog.
Analysis: distributions of the three classifications and applicable-pattern sets by story and system, retaining zero-positive stories. \\
\hline
RQ2 &
To what extent do the patterns classified as applicable by the LLM overlap with the frozen RAG baseline's final Top-5 candidates for the same requirement cases? &
Unit: one LLM-positive user-story--pattern pair, grouped by story.
Evidence: official applicability assessments and the frozen final Top-5 rankings.
Analysis: membership-intersection counts by story and overall, with the LLM-positive set as an explicit denominator; no applicability-accuracy metric. \\
\hline
RQ3 &
How do experts assess the applicability of the LLM-selected patterns and the adequacy, grounding, consequences, and perceived usefulness of the corresponding generated ADRs? &
Unit: one expert's rating of one selected pair and its ADR on one questionnaire item.
Evidence: consolidated expert responses to Q1--Q5 and optional comments for the prepared cases.
Analysis: item-specific ordinal distributions and case-level summaries with response counts; comments contextualize ratings. Aggregation and inclusion rules must be specified before consolidated analysis. \\
\hline
\end{tabular}
\end{table*}


RQ1 characterizes the model's classification outcomes rather than measuring
their correctness. RQ2 compares the patterns classified as applicable with a
separate retrieval baseline. This baseline retrieves catalog candidates and
reranks them to a final Top-5 without making an LLM-based applicability judgment.
Overlap is therefore a descriptive relationship between two output sets, not
evidence of either approach's precision, recall, or superiority.

ADR generation is a separate stage: each pair classified as applicable supplies
the requirement case, selected pattern, and applicability evidence for one
generated ADR. Generation completeness and stored provenance establish the
availability and traceability of these records, not their architectural quality.
RQ3 examines expert assessments of these selected pairs and their corresponding
ADRs, rather than adding another automatic applicability classification.

The prepared questionnaire measures decision adequacy for the stated
requirements (Q1), decision grounding (Q2), presentation of relevant
consequences and trade-offs (Q3), perceived usefulness as a decision record
(Q4), and pattern applicability (Q5). These are five separate expert-assessed
dimensions; they are not a validated composite quality score. Optional comments
provide context for the ratings. Neither implementation feasibility, cost
savings, novelty, nor actual improvement in architectural practice is directly
measured by these items.

Experts evaluate cases for a system they know, and only pairs selected by the
LLM are presented. Consequently, RQ3 can characterize expert judgments of the
presented recommendations but cannot identify missed applicable patterns or
estimate recall over the full set of assessed pairs. Expert agreement with a
recommendation does not by itself establish implementation success or objective
correctness. Repeated ratings from an expert and multiple cases from a system
must not be treated as independent replications when interpreting the results.
Conclusions remain bounded to the sampled requirements, catalog, generated
records, and participating experts.

% Claim-to-evidence matrix for Task 1.1. Paths are relative to the paper directory;
% these comments are provenance pointers, not external LaTeX build inputs.
% E1: ../20261012.0933-InitialArtifacts/adr4ai_app/experiments/pattern_selection/
%     assess_pattern.py and outputs/assessments.jsonl.
% E2: ../20261012.0933-InitialArtifacts/adr4ai_app/data/user_histories_orig.csv
%     and catalogs/patterns/.
% E3: ../20261012.0933-InitialArtifacts/adr4ai_app/experiments/rag/
%     README.md and outputs/rag_baseline.csv (not reproduction or pilot outputs).
% E4: ../20261012.0933-InitialArtifacts/adr4ai_app/research_results/
%     adr_generation_results.csv and compile_adr_results.py;
%     ../20261012.0933-InitialArtifacts/adr4ai_app/experiments/adr_generation/
%     generate_adr.py and outputs/*.json (excluding attempts/).
% E5: ../20261012.0933-InitialArtifacts/adr4ai_app/research_results/google_forms/
%     create_expert_evaluation_form.gs, EVALUATION_QUESTIONS, and README.md;
%     research_results/expert_evaluation_cases_pt.csv.
% E6: Future frozen/consolidated expert-response dataset and documented analysis.
%     Researcher reports preliminary responses; no response data inspected here.
%
% C1 / RQ1: generated label distributions and positive sets -> E1 + E2.
%     Allowed: descriptive counts/sets. Excluded: expert-confirmed correctness.
% C2 / RQ2: final Top-5 membership overlap -> E1 + E3, joined by story/pattern ID.
%     Allowed: intersection counts. Excluded: precision, recall, superiority.
% C3 / generation: one traceable ADR per selected pair -> E4 plus E1.
%     Allowed: completeness/provenance once verified. Excluded: semantic quality.
% C4 / RQ3 constructs: Q1--Q4 ADR dimensions; Q5 pattern applicability -> E5.
%     Allowed now: instrument description. Ratings/results require E6.
% C5 / RQ3 findings: observed expert ratings/comments -> E5 + E6.
%     Blocked pending consolidation; no synthetic or favorable mock findings.
%     Excluded: missed-pattern recall, population-wide effectiveness, deployment
%     success, cost savings, and a composite score without justification.
